\documentclass[10pt,a4paper]{amsart}
\usepackage[margin=19mm]{geometry}
\usepackage{amsmath,amssymb,mathtools}
\usepackage[scr=rsfs,cal=euler]{mathalfa}
\usepackage{xfrac}
\usepackage[unicode,hidelinks,bookmarks=false]{hyperref}
\newcommand{\ie}{i.e.\ }
\newcommand{\cf}{cf.\ }
\newcommand{\resp}{respectively\ }
\newcommand{\Subsection}{\subsection}
\providecommand{\nobreakdash}{\nobreak\hbox{-}}
\setlength{\emergencystretch}{2em}
\multlinegap0pt
\usepackage{bm}
\usepackage{bbm}
\usepackage{dsfont}
\usepackage{xspace}
\usepackage{cancel}
\usepackage{mathtools}
\usepackage{amsthm}
\makeatletter
\@ifundefined{theorem}{\newtheorem{theorem}{Theorem}}{}
\@ifundefined{claim}{\newtheorem{claim}[theorem]{Claim}}{}
\@ifundefined{definition}{\newtheorem{definition}[theorem]{Definition}}{}
\@ifundefined{lemma}{\newtheorem{lemma}[theorem]{Lemma}}{}
\makeatother
\DeclareMathOperator{\proj}{proj}
\newcommand{\B}{\mathcal{B}}
\newcommand{\CT}{\textmd{CT}}
\newcommand{\N}{\mathbb{N}}
\newcommand{\abs}[1]{\left\vert#1\right\vert}
\renewcommand{\ref}[1]{(\oldref{#1})}
\title[A Backward Ergodic Theorem for Uncountable-to-one Transforma]{A Backward Ergodic Theorem for Uncountable-to-one Transformations}
\author{Eric Wang}
\date{Source: arXiv:2605.29323v1 (CC BY 4.0)}
\begin{document}
\maketitle

\noindent\emph{Source and licence.} A Backward Ergodic Theorem for Uncountable-to-one Transformations. Version arXiv:2605.29323v1 (CC BY 4.0), distributed under the Creative Commons Attribution 4.0 International licence, as recorded in the arXiv licence metadata for this exact version and shown on the abstract page https://arxiv.org/abs/2605.29323v1. Full rights notice: licenses/arxiv-2605.29323-CC-BY-4.0.txt.\par\smallskip

\noindent\textit{Editorial scope.} Counted unit: the Theorem whose source label is \texttt{None} in the article (source0.tex lines 813--818, source label \texttt{None}), delivered together with the article's own complete original proof of it (source0.tex lines 819--898, one \texttt{proof} environment of 80 lines). No mathematics is written, repaired or re-derived: statement and proof are the article's own text line for line. Required context retained uncounted: JVN (lines 157--159); gluing (lines 459--464); localglobal (lines 701--706); The Theorem (lines 740--746); tiling (lines 750--752). Every definition, display and label of those lines is retained unchanged. Omitted from this package and not redistributed: 1--156, 160--458, 465--700, 707--739, 747--749, 753--812, 899--981. Nothing inside a retained excerpt is omitted or altered: every retained body is delivered exactly as the article's lines are written, so no omission receipt of any kind accompanies this package. Citations: the retained excerpts carry no citation command at all, so no bibliography entry of the article is transcribed and no reference is redistributed. Reference adaptation: none was needed and none was made; every reference inside the delivered unit resolves inside it, so no unresolved reference or citation is produced. Printed slips kept literal: no printed word, symbol, hypothesis or number of the counted statement or of its proof was altered. Layout and class: the article's own class is \texttt{article} with packages including geometry, amsmath, amsthm, amssymb, amsfonts, enumerate, marvosym, tipa, stmaryrd, mathtools, mathrsfs, gensymb, mathdots, url; the wrapper is the lane's pinned \texttt{amsart} editorial wrapper at 10pt on A4, which restates the theorem-like environments the retained text needs and adds the article's own macro definitions that the retained text needs (\texttt{\textbackslash B}, \texttt{\textbackslash CT}, \texttt{\textbackslash N}, \texttt{\textbackslash abs}, \texttt{\textbackslash proj}, \texttt{\textbackslash ref}), with extra wrapper packages (bm, bbm, dsfont, xspace, cancel, mathtools). The delivered numbering is the unit's own. Assets: no retained span contains a figure environment, a graphics-inclusion command, a PDF-page inclusion, a table or a TikZ picture; no image file is copied into this package, the package declares no asset dependency and no image file is redistributed with it. Licence: version arXiv:2605.29323v1 of this article is distributed under the Creative Commons Attribution 4.0 International licence, as recorded in the arXiv licence metadata for this exact version and shown on the abstract page https://arxiv.org/abs/2605.29323v1. A full rights notice is delivered at licenses/arxiv-2605.29323-CC-BY-4.0.txt. Characters: every retained excerpt is pure ASCII, so the delivered engine, XeLaTeX with its default Latin Modern text font, reports no missing character.
\begin{theorem}
    \label{JVN} (Jankov, von Neuman) Let $ X $ and $ Y $ be standard Borel. If $ P\subseteq X\times Y $ is analytic. Then there exists a $ \sigma(\Sigma^1_1) $-measurable uniformization of $ P $.
\end{theorem}

\begin{lemma}
    \label{gluing}(Gluing Lemma) Let $ x\in X $ and $ n\in \N $, and suppose $ A \subseteq T^{-n}\{x\} $ be measurable. Further suppose that for each $ y\in A $, that the function $ y\mapsto E_y $ is measurable as a function from $ X $ into $ \mathcal{Y} $ with $ E_y\subseteq T^{-1}\{y\} $. Then there exists a measurable set $ F\subseteq T^{-(n+1)}\{x\} $ so that for $ \mu_x^n $-a.e. $ y\in A $, $ \mu_y^1(F) = \mu_y^1(E_y) $. Furthermore, if $ f\in L^1 $, then 
    \[
    \int_{F} fd\mu_x^{n+1} = \int_A \int_{E_y} fd\mu_y^1d\mu_x^n(y)
    \].
\end{lemma}

\begin{lemma}\label{localglobal}
    Let $ N\in \N $, with $ f\in L^1(X,\B,\mu) $. Then
    \[
    \int_X fd\mu = \int \frac{1}{N+1}\sum_{i=0}^N \int_{T^{-i}\{x\}} f(y)d\mu^i_x(y) d\mu(x)
    \]
\end{lemma}

\begin{theorem}
    \label{The Theorem}(Backward Ergodic Theorem for arbitrary measure preserving transformations) Let $ X $ be a standard probability space, let $ T:X\rightarrow X $ be a measure-preserving transformation. Then for every $ f\in L^1(X) $, and for almost every $x \in X $, $ \abs{A_{f,E^x}} <\infty $ and the limit
    \[
    \lim_{m_x(E_x)\nearrow \infty} A_{f,E_x} 
    \]
    exists and equals $ \overline{f}= \mathbb{E}(f|\mathcal{B}_{\mathcal{I}}) $, the conditional expectation of $ f $ with respect to the $ \sigma $-algebra of $ T $-invariant subsets, where $ E_x$ varies over the set of coherent trees rooted at $x $.
\end{theorem}

    \begin{claim}\label{tiling}
        For any $ \varepsilon>0 $, there exists $ N\in \N $ and subset $ X'\subset X $ of measure $ \geq 1-\varepsilon $ such that for all $ x\in X' $, there exists a coherent subtree $ S $ of complete tree $ \triangleright_T^N\cdot x $ so that $ m_x(S')\geq (1-\varepsilon)m_x(\triangleright_T^N\cdot x) $ and with $ A_{f-g,S} >0 $. 
    \end{claim}

\begin{theorem}
    (Backward Ergodic Maximal Theorem) Let $ T $ be as in \ref{The Theorem}. Then for every $ f\in L^1(X,\mu) $, if we define $ f^\star  = \sup_{E_x\in \CT_x} A_{f,E_x} $ for each $ x\in X $. Then for any $ \lambda\in \mathbb{R} $, 
    \[
    \int_{\{f^\star>\lambda\}} f d\mu \geq \lambda \mu(\{f^\star>\lambda\})
    \]
\end{theorem}

\begin{proof}
    First, we see that $ f^\star $ is $ \sigma(\Sigma^1_1) $-measurable. To see this, let $ \alpha\in \mathbb{R} $. We will show that $ \{f^\star>\alpha\} $ is analytic. Consider $ \mathcal{A}\subseteq X\times \mathcal{Y} $ defined as follows: $ (x,E)\in \mathcal{A} $ if and only if:
    \begin{enumerate}
        \item $ E\in \CT_x $
        \item $ A_{f,E}>\alpha $
    \end{enumerate}
    As we have seen earlier, the two conditions defining $ \mathcal{A} $ are Borel. Hence, $ \{f^\star>\alpha\} = \proj_X(\mathcal{A}) $, and is thus analytic.\\\\
    Now, let $ \lambda\in \mathbb{R} $. Let $ Y = \{x\in X:f^\star>\lambda\} $. We want to show that
    \[
    \int_Y fd\mu \geq \lambda \mu(Y) 
    \]
    It suffices to show that for every $ \varepsilon>0 $, $\int_Y(f-\lambda)d\mu \geq -\varepsilon $.
    \begin{definition}\label{def for min}
        Let $ x\in Y $. Then let $ E_x $ be a coherent tree that witnesses $ x\in Y $. We say $ E_x $ is \textbf{minimal} if there is no proper \textbf{subtree}, $ F_x $ such that $ A_{f,F_x} > \lambda $. \\\\
        More precisely, $ F_x $ is a \textbf{subtree} of $ E_x $ if for each $ i $,
        \[
        \mu^i_x(F_x^i\setminus E_x^i) = 0
        \]
        And, if $ F_x $ is a subtree of $ E_x $ such that $ A_{f,F_x}>\lambda $, then $ E_x $ and $ F_x $ agree ($\mu_x^i$)-almost-everywhere at each level. In particular, $ m_x(F_x) = m_x(E_x) $.
    \end{definition}
    \begin{claim}\label{minimal gives right thing}
        Suppose $ x\in Y $, and $ E_x $ is a minimal witness as in \ref{def for min}. Then for $ m_x $-almost-every $ y\in E_x $, $ y\in Y $. More precisely, for each $ 0\leq i \leq \ell(E_x) $, 
        \[
        \mu_x^i(E_x^i\setminus Y) = 0
        \]
    \end{claim}
    \noindent We call the property of $ E_x $ in the hypothesis of \ref{minimal gives right thing}, \textbf{$Y$-saturated}. 
    \begin{proof}
        (of claim) Suppose, for the sake of contradiction, that for some $ 0\leq i \leq \ell(E_x) $, $ \mu^i_x(E_x^i\setminus Y) > 0 $. Then we can simply remove the portion of $ E_x $ behind that subset and obtain a proper subtree of strictly smaller mass, on which the average of $ f $ is still greater than $ \lambda $, contradicting minimality of $ E_x $.
    \end{proof}
    \begin{claim}
        For every $ x\in Y $, there exists a minimal witnessing tree $ E_x $.
    \end{claim}
    \begin{proof}
        (of claim) Let $ x\in Y $. Since $ \sup_{E\in \CT_x} A_{f,E} > \lambda $, let $ E\in \CT  $ be a witness. We perform a transfinite induction on countable ordinals:
        \begin{itemize}
            \item Let $ \mathcal{E}_0 = E $.
            \item Suppose $ i<\omega_1 $, and we have $\mathcal{E}_i \in \CT_x $ witnessing $ A_{f,\mathcal{E}} >\lambda $. Let $ \mathcal{E}_{i+1} $ be a proper subtree of $ \mathcal{E}_i $ with $ m_x(\mathcal{E}_{i+1})\leq m_x(\mathcal{E}_i) $. 
            \item If $ i<\omega_1 $ is a limit ordinal, let $ \mathcal{E}_i = \bigcap_{j<i}\mathcal{E}_i $ (where we do coordinate-wise intersection). Then $ \mathcal{E}_i \in \CT_x $, and $ m_x(\mathcal{E}_i)\leq m_x(\mathcal{E}_j) $ for each $ j<i $.
        \end{itemize}
        Then the function $ \zeta:\omega_1\to [1,\infty) $, $ i\mapsto m_x(\mathcal{E}_i) $ must stabilize at some countable ordinal $ i $. Then $ \mathcal{E}_i $ is the minimal witness.
    \end{proof}
    \begin{claim}\label{maximaluniformization}
        There is a $ \sigma(\Sigma^1_1) $-measurable function $ \gamma:X\to \mathcal{Y} $ such that for every $ x\in Y $, $ \gamma(x) $ is $ Y $-saturated. 
    \end{claim}
    \begin{proof}
        (of claim) Define $ \mathcal{A}\subseteq X\times \mathcal{Y} $ where $ (x,E)\in \mathcal{A} $ if and only if the following are true:
        \begin{enumerate}
            \item $ E\in \CT_x $
            \item $ A_{f,E}>\lambda $
            \item $ E $ is $ Y $-saturated. 
        \end{enumerate}
        We have already established in earlier proofs that conditions (1) and (2) are Borel. And it's easy to see that condition (3) is also Borel.\\\\
        Then by Jankov von-Neumann, \ref{JVN}, there exists a $ \sigma(\Sigma^1_1) $-uniformization $ \gamma:X\to \mathcal{Y} $, as desired.
    \end{proof}
    \noindent With the uniformization, we can proceed with a proof similar to that in \ref{The Theorem}. Let $ \varepsilon>0 $. As in \ref{The Theorem}, we can assume that $ f $ is bounded below. 
    \begin{claim}
        For any $ \varepsilon>0 $, there exists $ N\in \mathbb{N} $ and subset $ X'\subset X $ of measure $ \geq 1-\varepsilon $ such that for all $ x\in X' $, there exists a coherent subtree $ S $ of complete tree $ \triangleright^N_T\cdot x $ so that $ m_x(S)\geq (1-\varepsilon)m_x(\triangleright_T^N\cdot x) $ and with
        \[
        A_{f,S}>\lambda 
        \]
    \end{claim}
    \begin{proof}
        The proof of the claim follows identically to that of \ref{tiling} in which we use a greedy argument using the gluing lemma, \ref{gluing}, with the measurable uniformization obtained in \ref{maximaluniformization}.
    \end{proof}
    \noindent We further have, using the same argument as earlier,
    \begin{claim}
        For each $ x\in X' $ above,
        \[
        A_{1_Y\cdot(f-\lambda),\triangleright^N_T\cdot x} \geq -\frac{\varepsilon}{2}
        \]
    \end{claim}
    \noindent Finally, by \ref{localglobal}, we see
    \begin{align*}
        \int_Y (f-\lambda)d\mu &= \int A_{1_Y\cdot (f-\lambda),\triangleright_T^N\cdot x}d\mu\\
        &= \int_{X'} A_{1_Y\cdot (f-\lambda),\triangleright^N_T\cdot x} d\mu + \int_{X\setminus X'}A_{1_Y\cdot (f-\lambda),\triangleright^N_T\cdot x} d\mu \\
        &\geq \frac{\varepsilon}{2} + \int_{X\setminus X'} A_{1_Y\cdot (f-\lambda),\triangleright^N_T\cdot x}d\mu
    \end{align*}
    Taking $ X' $ to be sufficiently close measure to $ X $ and since $ f $ can be assumed to be bounded below, we can conclude $ \int_{Y}(f-\lambda)d\mu \geq -\varepsilon $, as desired.
\end{proof}

\end{document}
